\documentclass{article}
\usepackage[left=2cm, right=2cm, top=2cm, bottom=2cm]{geometry}
\usepackage{graphicx}
\usepackage{amsmath}
\usepackage{amssymb}
\usepackage{amsthm}
\usepackage{fancyhdr}
\usepackage{verbatim}
\usepackage{listings}
\usepackage{xcolor}
\usepackage{pdfpages}
\usepackage{cancel}
\usepackage{physics}
\usepackage{esint}

\lstset{
    language=C++,
    basicstyle=\ttfamily\footnotesize,
    keywordstyle=\color{blue},
    commentstyle=\color{green},
    stringstyle=\color{red},
    numbers=left,
    numberstyle=\tiny\color{gray},
    frame=single,
    breaklines=true,
    captionpos=b,
}


\title{MTH 553: Lecture \# 37}
\author{Cliff Sun}

\newtheorem{theorem}{Theorem}[section]
\newtheorem{lemma}[theorem]{Lemma}
\newtheorem{definition}[theorem]{Definition}
\newtheorem{conjecture}[theorem]{Conjecture}
\newtheorem{proposition}[theorem]{Proposition}
\newtheorem{corollary}[theorem]{Corollary}
\newtheorem{one minute paper}[theorem]{One Minute Paper}

\pagestyle{fancy}
\lhead{\textbf{\thepage}\ \ \nouppercase{\rightmark}}
\chead{MTH 553: Lecture \# 37}
\rhead{Cliff Sun}

\begin{document}

\maketitle

\section*{Lecture Span}
\begin{enumerate}
    \item Similarity Method
\end{enumerate}

\section*{Similarity Method}

A lot of PDEs have \underline{similarity} solutions defined as 

\begin{equation}
    u(x,t) = t^{\alpha}w(t^{\beta} x)
\end{equation}

For some constants $\alpha, \beta$. For instance, for the one dimensional heat equation 

\begin{equation}
    u_t = u_{xx}
\end{equation}

Then we can substitute this similarity solution into this PDE to obtain 

\begin{equation}
    u_t = t^{\alpha - 1}\left( \alpha w(t^{\beta}x) + \beta(t^{\beta}x)w'(t^{\beta}x) \right)
\end{equation}
\begin{equation}
    u_{xx} = t^{\alpha + 2\beta}w''(t^{\beta}x)
\end{equation}

To get the $t$ terms to disappear, let $\alpha + 2\beta = \alpha -1$ implies that 

\begin{equation}
    \beta = -\frac{1}{2}
\end{equation}

Then, we can obtain an equation in terms of $w$, 

\begin{equation}
    w''(y) = \alpha w(y) + \beta yw'(y)
\end{equation}
\begin{equation}
    w''(y) = \alpha w(y) - \frac{1}{2} yw'(y)
\end{equation}

Where $y = t^{\beta} x \in \mathbb{R}$. In essence, we have reduced the PDE into an ODE. T
hen, we can solve this ODE for $w$. Let $\alpha = \beta$ (note, this choice permits conservation of heat). Therefore, $\alpha = \beta = -1/2$. Then, 

\begin{equation}
    \int_{-\infty}^{\infty}t^{\beta}w(t^{\beta}x)dx = \int_{-\infty}^{\infty}w(y)dy
\end{equation}

Note that this is independent of $t$. On the LHS, this is the total heat (since you are integrating the heat equation over $x$). And this doesn't depend on $t$ which means that 
$\alpha = \beta = -1/2$ permits heat conservation. We solve 

\begin{align}
    w'' &= -\frac{1}{2}\left( w + yw' \right)\\
    w'' &= -\frac{1}{2}\left( yw \right)' \\
    w' &= -\frac{1}{2}yw + c_1, \quad \text{choose $c_1 = 0$} \\
    w(y) &= C\exp(-\frac{y^2}{4})
\end{align}

Gaussian equation. Therefore,

\begin{align}
    u(x,t) &= t^{-1/2}w(t^{-1/2}x)\\
    &= \frac{C}{(4\pi t)^{1/2}}\exp(-x^2/4t) 
\end{align}

\subsection*{Remarks}

\begin{enumerate}
    \item The choice of $\beta$ was forced to eliminate the leading powers of $t$
    \item Other choices of $\alpha$ are mathematically possible. Although not physically possible. 
\end{enumerate}

\section*{Traveling wave for a non-linear PDE}

This is called the Korteweg-de Vries equation (KdV). 

\begin{equation}
    u_t + 6uu_{x} + u_{xxx} = 0
\end{equation}

Note that this equation is non-linear and third order. Interpretation: $u$ is the height of water or electrical wave in one dimension.
Let's try a traveling wave. 

\begin{equation}
    u(x,t) = v(x - ct)
\end{equation}

We then obtain 

\begin{equation}
    -cv'(y) + 6v(y)v'(y) + v'''(y) = 0
\end{equation}

Where $y = x - ct$. We can integrate everything 

\begin{equation}
    -c v + 3v^2 + v'' = a
\end{equation}

Where $a$ is a constant. Multiply by $2v'$ and integrate 

\begin{equation}
    (v')^2 = cv^2 - 2v^3 + 2av + b
\end{equation}

Where $b$ is a constant. This is a non-linear ODE that is first order. Suppose $v,v',v''$ tend to $0$ as $y \rightarrow \pm \infty$. This is a solitary wave. 
Thus, $b=0$ and $a=0$. Therefore, for a solitary wave, we obtain 

\begin{equation}
    v' = \pm v\sqrt{c - 2v}
\end{equation}

This is a seperable non-linear ODE. Change variable, let $$v(y) = \frac{c}{2}\sech^2(\theta(y))$$ 

Then we obtain:

\begin{align}
    \pm v\sqrt{c - 2v} &= \pm \frac{c}{2}\sech^2(\theta)\sqrt{c - c\sech^2\theta} \\
    &= \pm\frac{c^{3/2}}{2}\sech^2\theta \tanh(\theta)  
\end{align}
\begin{equation}
    v'(y) = -c\sech^2\theta \tanh\theta
\end{equation}

Thus 

\begin{equation}
    \frac{d\theta}{dy} = \pm \frac{\sqrt{c}}{2} \implies \theta = \pm \frac{\sqrt{c}}{2}\left( y - \phi \right)
\end{equation}

Thus, 

\begin{equation}
    \boxed{v(y) = \frac{c}{2}\sech^2\left( \frac{\sqrt{c}}{2}[y - \phi] \right)}
\end{equation}

This is our solitary wave with height $c/2$ with speed $c$. They are related! In particular, tall waves travel faster. When tall soliton 
wave overtakes small wave, then the tall wave gets a bump forwards and the small wave gets pushed back. But they keep the same shape for some reason. 

A traveling wave is also a similarity solution in logarithmic coordinates. 

\begin{equation}
    v(x - ct) = v\left( \log \frac{\xi}{\tau^c} \right) 
\end{equation}

Where $\xi = \exp(x)$ and $\tau = e^{t}$. 

\end{document}